\section{Conclusions}

We have proposed a novel information-theoretic approach for Bayesian
optimization. Our method, predictive entropy search (PES), greedily maximizes
the amount of one-step information on the location $\xopt$ of the global maximum
using its posterior differential entropy.  Since this objective function is
intractable, PES approximates the original objective using a reparameterization
that measures entropy in the posterior predictive distribution of the function
evaluations.  PES produces more accurate approximations than Entropy Search
(ES), a method based on the original, non-transformed acquisition function.
Furthermore, PES can easily marginalize its approximation with respect to the
posterior distribution of its hyper-parameters, while ES cannot.  Experiments
with synthetic and real-world functions show that PES often outperforms ES in
terms of immediate regret. In these experiments, we also observe that PES often
produces better results than expected improvement (EI), a popular heuristic for
Bayesian optimization.  EI often seems to make excessively greedy decisions,
while PES tends to explore more. As a result, EI seems to perform better for
simple objective functions while often getting stuck with noisier objectives or
for functions with many modes.
